
\subsection{Deviatoric Basis}\label{deviatoric-basis}

Define 
\[
\mathbf{D}:=\operatorname{dev}_{\Omega}\mathbf{I}=\mathbf{I}-\tfrac12(\mathrm{tr}\,\mathbf{I})\,\mathbf{1},\qquad
N:=\sqrt{\mathbf{D}\!:\!\mathbf{D}}\quad(\text{so }N>0\text{ unless the section is in-plane isotropic}).
\] 
Then define the orthonormal, symmetric deviatoric basis tied to \(\mathbf{I}\) and
\([\mathbf{i}^{\times}]\): 
$$
\mathbf{U}:=\frac{D}{2 D:D}, \quad \mathbf{~B}:=\frac{\operatorname{sym}\left(\left[\mathbf{i}^{\times}\right] D\right)}{D: D}=\frac{\left[\mathbf{i}^{\times}\right] D-D\left[\mathbf{i}^{\times}\right]}{D: D} .
$$
and
\[
\overline{\mathbf{U}}=\frac{\operatorname{dev}_{\Omega}\mathbf{I}}{\sqrt{\mathbf{D}\!:\!\mathbf{D}}},\qquad
\overline{\mathbf{V}}=\frac{\mathrm{sym}([\mathbf{i}^{\times}]\operatorname{dev}_{\Omega}\mathbf{I})}{\sqrt{\mathbf{D}\!:\!\mathbf{D}}}
\]
or
\[
\overline{\mathbf{U}}:=\frac{\mathbf{D}}{N},\qquad
\overline{\mathbf{V}}:=\frac{\mathrm{sym}\big([\mathbf{i}^{\times}]\,\mathbf{D}\big)}{N}
=\frac{[\mathbf{i}^{\times}]\,\mathbf{D}-\mathbf{D}\,[\mathbf{i}^{\times}]}{2N}
\]
with the properties 
\(\overline{\mathbf{U}}\!:\!\overline{\mathbf{U}}=\overline{\mathbf{V}}\!:\!\overline{\mathbf{V}}=1\)
and \(\overline{\mathbf{U}}\!:\!\overline{\mathbf{V}}=0\). 
Any section deviator \(\operatorname{dev}_{\Omega}\bm{R}\) decomposes uniquely as
\[
\operatorname{dev}_{\Omega}\bm{R}=\underbrace{(\operatorname{dev}_{\Omega}\bm{R}:{\mathbf{U}})}_{a}\,{\mathbf{U}}
+\underbrace{(\operatorname{dev}_{\Omega}\bm{R}:{\mathbf{V}})}_{b}\,{\mathbf{V}}.
\]
In
principal axes these reduce to:
\begin{equation}
\overline{\mathbf{U}}=\frac{1}{\sqrt{2}}\left[\begin{array}{cc}
1 & 0 \\
0 & -1
\end{array}\right], \quad \overline{\mathbf{V}}=\frac{1}{\sqrt{2}}\left[\begin{array}{ll}
0 & 1 \\
1 & 0
\end{array}\right]
\end{equation}


In a principal system, \(\mathbf{D}\!:\!\mathbf{D}\) expands as \[
\mathbf{D}\!:\!\mathbf{D}
=\mathbf{I}\!:\!\mathbf{I}-\tfrac12(\mathrm{tr}\,\mathbf{I})^2
=\tfrac12(\mathrm{tr}\,\mathbf{I})^2-2\,\det\mathbf{I}
=\tfrac12(I_y-I_z)^2+2\,I_{yz}^2.
\]

In a basis \({\mathbf{e}_1,\mathbf{e}_2}\) that diagonalizes
\(\mathbf{I}\), these are $a=\tfrac12(y^2-z^2)$, $b=yz$. 



\begin{itemize}
    \item 
Invariant formula (coordinate-free):

$$
\mathbf{D}: \mathbf{D}=\mathbf{I}: \mathbf{I}-\frac{1}{2}(\operatorname{tr} \mathbf{I})^2=\frac{1}{2}(\operatorname{tr} \mathbf{I})^2-2 \operatorname{det} \mathbf{I}
$$
where the two right-hand expressions are equal in 2D by Cayley-Hamilton: $\mathbf{t r I}^{\mathbf{2}}=(\mathbf{t r I})^{\mathbf{2}}-\mathbf{2} \operatorname{det} \mathbf{I}$.
\item Component form (with $\mathbf{I}=\left[\begin{array}{cc}I_y & -I_{y z} \\ -I_{y z} & I_z\end{array}\right]$ ):

$$
\mathbf{D}: \mathbf{D}=\frac{1}{2}\left(I_y-I_z\right)^2+2 I_{y z}^2=\frac{1}{2}\left(I_y^2+I_z^2\right)-I_y I_z+2 I_{y z}^2
$$

\end{itemize}

So $\mathbf{D}: \mathbf{D} \geq \mathbf{0}$, and it vanishes iff the section inertia is isotropic in the plane ( $\mathbf{I} \propto \mathbf{1}$, i.e., equal principal moments and zero product of inertia).

Let $\{\mathbf{e}_{1}, \mathbf{e}_{2}\}$ be the orthonormal eigenvectors of $\mathbf{I}$ in the section plane, with $\mathbf{e}_{1} \times \mathbf{e}_{2}=\mathbf{i}$. 
Define the two constant deviatoric tensors tied to $\mathbf{I}$:

$$
\mathbf{U}_I:=\mathbf{e}_{1} \otimes \mathbf{e}_{1}-\mathbf{e}_{2} \otimes \mathbf{e}_{2}, 
\quad
\mathbf{V}_I:=\frac{1}{2}(\mathbf{e}_{1} \otimes \mathbf{e}_{2}+\mathbf{e}_{2} \otimes \mathbf{e}_{1}) .
$$
or
\[
\mathbf{U}=\mathrm{diag}(1,-1),\qquad
\mathbf{V}=\begin{bmatrix}
0&1\\
1&0
\end{bmatrix}
\]



\subsection{Cubic vector fields}\label{cubic-vector-fields}

\subsubsection{The Poisson cubic}\label{the-poisson-cubic}

A clean expression is just the \(\overline{\mathbf{V}}\) projection twice: 
\[
\mathbf{f}(\bm{r})\;=\;b\;{\mathbf{V}}\,\bm{r}
\;=\;\bigl(\operatorname{dev}_{\Omega}\bm{R}\,\cdot\,{\mathbf{V}}\bigr)\ \bigl({\mathbf{V}}\,\bm{r}\bigr)
\] 
In principal axes: \(b=yz\) and \(\overline{\mathbf{V}}\bm{r}=(z,y)\),
hence \(\mathbf{f}=yz\,(z,y)=(y z^2,\ z y^2)\). PDE identity:
\(\Delta \mathbf{f}=2\,\bm{r}\) (so for any load vector \(\bm{p}\),
\(\Delta(\mathbf{f}\cdot\bm{p})=2\,\bm{r}\cdot\bm{p})\).

A couple of useful equivalent forms: 
\[
\mathbf{f} = b\,{\mathbf{V}}\,\bm{r}
= b\,\big(\operatorname{sym}([\mathbf{i}^{\times}]{\mathbf{U}})\big)\bm{r}
\quad\text{since}
\quad {\mathbf{V}}=\operatorname{sym} \bigl([\mathbf{i}^{\times}]{\mathbf{U}}\bigr).
\]

\subsubsection{The harmonic cubic}\label{the-harmonic-cubic}

This one is a linear combination of the ``radial'' cubic and the
\({\mathbf{V}}\) channel: 
\[
\mathbf{h}(\bm{r})\;=\;2\,a\,\bm{r}\;-\;4\,b\;{\mathbf{V}}\,\bm{r}
\;=\;2\,(\operatorname{dev}_{\Omega}\bm{R}:{\mathbf{U}})\,\bm{r}\;-\;4\,(\operatorname{dev}_{\Omega}\bm{R}:{\mathbf{V}})\,({\mathbf{V}}\,\bm{r})\ 
\] 
In principal axes: \(a=\tfrac12(y^2-z^2)\), \(b=yz\),
\({\mathbf{V}}\bm{r}=(z,y)\), so 
\[
\mathbf{h}_y= y^3-3yz^2,\quad
\mathbf{h}_z= z^3-3zy^2,
\]
and \(\Delta \mathbf{h}=\mathbf{0}\) (harmonic). 
\[
\mathbf{h}(\bm{r})=(0,\;y^3-3y z^2,\;z^3-3z y^2)
\]

\subsubsection{Isotropic Cubic}\label{isotropic-cubic}

The matrix \[
Q[\bm{r}^\times] =
\begin{bmatrix}
0&0&0\\
y^{2} z + z^{3} & 0 & 0\\
* y^{3} - y z^{2} & 0 & 0
\end{bmatrix}
\] 
has first column
\(\bm{Q}(\mathbf{i}\times\bm{r}) = (\bm{r}\cdot\bm{r})\,\mathbf{i}\times\bm{r}\).
That gives the \textbf{``isotropic'' cubic swirl}
\((\bm{r}\cdot\bm{r})\;\mathbf{i}\times\bm{r}\), whereas
\(\mathbf{h}\) is the \textbf{anisotropic harmonic} combination above.


\subsubsection{Summary}\label{summary}

\[
\boxed{\ \mathbf{f}=b\,{\mathbf{V}}\,\bm{r},\qquad
\mathbf{h}=2\,a\,\bm{r}-4\,b\,{\mathbf{V}}\,\bm{r}\ ,\quad
a=\operatorname{dev}_{\Omega}\bm{R}:{\mathbf{U}},\ \ b=\operatorname{dev}_{\Omega}\bm{R}:{\mathbf{V}}\ .\ }
\]




\subsection{3) Love's Neumann data}\label{loves-neumann-data}

Start from your original bracket and substitute
\(\mathbf{Q}=(\bm{r}\cdot\bm{r})\mathbf{1}-\bm{R}\),
\(\mathbf{S}=2\,\mathrm{Off}(\operatorname{dev}_{\Omega}\bm{R})\), then collect
isotropic and deviatoric terms. You get the compact tensor: \[
\boxed{
\mathbf{A}(\bm{r})
=\tfrac{1}{2}(\bm{r}\cdot\bm{r})\,\mathbf{1}
+\ (\nu-1)\,a\,{\mathbf{U}}
+\ (2\nu+1)\,b\,{\mathbf{V}}}
\] with
\(a=\operatorname{dev}_{\Omega}\bm{R}:{\mathbf{U}}\) and \(b=\operatorname{dev}_{\Omega}\bm{R}:{\mathbf{V}}\)
as above.

Therefore Love's BVP can be written as \[
\boxed{
\left\{
\begin{aligned}
\Delta\,\boldsymbol{\chi}&=0 &&\text{in }\bar\Omega,\\[2pt]
\partial_{\bm{\nu}}\boldsymbol{\chi}&= -\,\mathbf{A}(\bm{r})\,\bm{\nu}
&&\text{on }\partial\bar\Omega,\\[2pt]
\displaystyle\int_{\Omega}\boldsymbol{\chi}\,\mathrm{d}A&=\mathbf{0},\qquad
\int_{\Omega}\bm{r}\,\mathrm{d} A=\mathbf{0}
\end{aligned}
\right.}
\]

\paragraph{Remarks}\label{two-equivalent-minimalist-rewrites-you-might-like}

\begin{itemize}
\tightlist
\item
  Using only \(\operatorname{dev}_{\Omega}\bm{R}\) and the projections onto
  \(\overline{\mathbf{U}}\) and \(\overline{\mathbf{V}}\): \[
  \mathbf{A}
  =\frac{(\bm{r}\cdot\bm{r})}{2}\mathbf{1}
  +(\nu-1)\,(\operatorname{dev}_{\Omega}\bm{R}\!:\!\overline{\mathbf{U}})\,\overline{\mathbf{U}}
  +(2\nu+1)\,(\operatorname{dev}_{\Omega}\bm{R}\!:\!\overline{\mathbf{V}})\,\overline{\mathbf{V}}.
  \]
\item
  Using a split into diagonal/off-diagonal parts: \[
  \mathbf{A}
  =\frac{(\bm{r}\cdot\bm{r})}{2}\,\mathbf{1}
  +(\nu-1)\,\mathrm{Diag}(\operatorname{dev}_{\Omega}\bm{R})
  +(2\nu+1)\,\mathrm{Off}(\operatorname{dev}_{\Omega}\bm{R}).
  \]

\item In a \textbf{Principal-axis frame} with \({\mathbf{U}}=\mathrm{diag}(1,-1)\), \({\mathbf{V}}\)
off-diag, and \(a=\tfrac12(y^2-z^2),\,b=yz\), \[
\mathbf{A}
=\frac{y^2+z^2}{2}\,\mathbf{1}
+(\nu-1)\,\frac{y^2-z^2}{2}\,\mathrm{diag}(1,-1)
+(2\nu+1)\,yz
\begin{bmatrix}
0&1 \\
1&0\end{bmatrix}
\]
and 
\(-\mathbf{A}\,\bm{\nu}\)
reproduces exactly
\[
-\left[\tfrac{\nu}{2}\bm{R}+\Big(1-\tfrac{\nu}{2}\Big)\mathbf{Q}+\tfrac{2+\nu}{2}\mathbf{S}\right]\bm{\nu}
\] 
i.e.~Love's original boundary traction, term by term.
\end{itemize}

